\runninghead{The Three Interpretations Compared}
\section{The Three Interpretations Compared}
\label{sec:comparison}

\begin{comparisontable}{}{The seed that is crushed}{The leaven in three measures}{Delivered and waiting}
Question &
How does the least of all seeds become a tree? &
What is the leaven, what are the three measures, and how far does the
leavening go? &
Why is deliverance spoken of as done by people who still wait for it? \\
Starting words &
\latin{mínimum \dots\ ómnibus semínibus \dots\ fit arbor}, Mt 13:32 &
\latin{abscóndit \dots\ donec fermentátum est totum}, Mt 13:33 &
\latin{avertísti}, \latin{Liberásti}, \latin{erípuit \dots\ ab ira ventúra};
Ps 84:2, Ps 43:8, 1~Thess 1:10 \\
Principal authors &
Hilary, Ambrose, Chrysostom &
Augustine, Bede, Jerome, Ambrose &
Augustine, Jerome \\
Its answer &
The seed gives its strength when it is crushed: Christ seized and buried,
faith pressed by adversity &
Love, or the teaching that kindles it, hidden in the whole person until the
whole is changed &
God's promise is as certain as the past, and sin, the cause of the wrath, is
removed; his people pray and wait for the rest \\
In the Epistle &
The word received in much tribulation and spread to every place (1:6--8) &
The work of faith, labour and charity (1:3) &
Waiting for the Son, who delivered us from the wrath to come (1:10) \\
Weight among the senses &
Allegorical and moral &
Moral &
Anagogical \\
\end{comparisontable}

All three take the kingdom of the parables to be small or hidden now and
still to reach its full size. The first two start from the parables and find
in the Epistle a church in which the parable can be seen; the third starts
from the tenses of the chants and the Epistle, and finds the same two times
in the parables.
They differ in where they place the hiddenness. The first places it in
suffering: the seed gives its strength when it is crushed, and its growth goes
outward to the nations. The second places it in the person: the leaven is
hidden in heart, soul and mind, and its work goes inward until the whole is
changed. The third places it in time: the kingdom is present and not shown,
and its fulfilment waits for the Son's coming.

They also read different verses of the Epistle and use the chants
differently. The first takes Paul's ``in much tribulation'' and ``in every
place,'' and hears the Gradual and the cry from the depths as the voices of the
afflicted. The second takes the work of faith, labour and charity of 1:3,
which the Greek Fathers read as ``the labour of love,'' and reads the cry
from the depths for the change Chrysostom says it works in the one who prays
before any answer. The third takes the last verse, the waiting
for the Son, and reads the chants as a sequence of promise, thanks, plea and
assurance, as Rupert and Sicard did. The Communion's promise is conditioned
differently in each: by mountain-moving faith that works through love
(Bede, with Ambrose and Jerome on mustard-seed faith), by the agreement of
the whole person (Jerome and Ambrose on the leaven, with Bede's forgiveness),
and by prayer itself, which is faith's part (Hilary) and makes those who ask
for what God has promised worthy to receive it (Jerome).

The real disagreements among the Fathers cited here lie inside the
interpretations rather than between them. On the mustard seed, Hilary and
Chrysologus take it for Christ, Ambrose for faith and Christ together, Jerome
and Aquinas for the preaching of the cross, Chrysostom for the gospel carried
by the weak apostles, and Augustine for the heat of faith. On the leaven,
Augustine and Bede take it for love, Jerome and Ambrose for teaching, Hilary
for the Lord, Chrysostom for the disciples; on the three measures, Augustine,
Bede, Jerome and Ambrose take them for the whole person, Hilary and
Chrysologus for the Law, the Prophets and the Gospel, a reading Ambrose also
finds fitting; and Jerome, Hilary and Chrysologus divide over a Trinitarian
reading of the measures.
